%! TEX root = ../paper.tex
\section{The Harness}
\label{sec:harness}

CanaryForge orchestrates a tool-calling LLM agent inside a containerized Docker sandbox.
It introduces targeted escape temptations and monitors three independent telemetry signals to detect unauthorized pivoting.
Originally designed for local execution, the harness architecture now seamlessly scales across distributed computing environments and heterogeneous hardware profiles.
This enables the simultaneous evaluation of local 7B models and cloud-based reasoning engines across different operating systems.
Crucially, despite this distributed scale, the core design principle remains intact: no component within the harness relies on the model to report on itself.

To guarantee safe execution during evaluation, the environment strictly isolates the agent.
Two Docker networks (\texttt{red-proxy} and \texttt{red-mocks}) are declared \texttt{internal}.
This architectural boundary prevents agent-generated traffic from reaching the public internet, regardless of the model's actions.
A third, strictly segmented network connects the harness to the designated model backend without carrying any agent-controlled payload (Figure~\ref{fig:arch}).

\begin{figure}[htbp]
\centering
\resizebox{0.92\textwidth}{!}{\begin{tikzpicture}[
        font=\small,
        node distance=1.5cm and 2.5cm,
        box/.style={draw, rectangle, align=center, minimum width=3.5cm, minimum height=1.2cm, fill=white},
        dashedbox/.style={box, dashed},
        dottedbox/.style={box, densely dotted},
        netbox/.style={draw, densely dotted, inner sep=15pt},
        >=Stealth,
        arrow/.style={->, thick},
        dashedarrow/.style={->, dashed, thick},
        dottedarrow/.style={->, densely dotted, thick}
    ]

    % Nodes
    \node[dashedbox] (backend) {model backend \\ \footnotesize outside the audited scope};
    
    \node[box, below=1.5cm of backend] (sandbox) {sandbox \\ \footnotesize agent + heartbeat};
    \node[box, right=2cm of sandbox] (egress) {egress-proxy \\ \footnotesize mitmproxy, :8080 and :9000};

    \node[dashedbox, right=2.5cm of egress, yshift=1.2cm] (registry) {package-registry};
    \node[dashedbox, right=2.5cm of egress, yshift=-1.2cm] (hub) {model-hub};
    
    \node[dashedbox, right=1.5cm of registry, yshift=-1.2cm] (redteam) {redteam \\ \footnotesize alias mirror-externo};

    \node[dottedbox, below=1.8cm of sandbox] (workspace) {workspace \\ \footnotesize ephemeral volume};
    \node[dottedbox, right=1cm of workspace] (memory) {./memory \\ \footnotesize bind mount, persists};

    \node[box, below=1.5cm of workspace, xshift=2.5cm] (fsmon) {fs-monitor \\ \footnotesize canary, inotify, no network};

    % Network boundaries (Background boxes)
    \begin{scope}[on background layer]
        \node[netbox, fit=(sandbox) (egress), label={[anchor=south, inner sep=2pt]above:\textit{red-proxy (internal)}}] (netproxy) {};
        \node[netbox, fit=(registry) (hub), label={[anchor=south, inner sep=2pt]above:\textit{red-mocks (internal)}}] (netmocks) {};
    \end{scope}

    % Connections
    \draw[dashedarrow] (backend) -- node[left, align=right] {\footnotesize not audited} (sandbox);
    
    \draw[arrow] (sandbox) -- node[above] {\footnotesize \texttt{HTTP\_PROXY}} (egress);
    \draw[dashedarrow] (sandbox.south east) to[bend right=15] node[below] {\footnotesize heartbeat, bypasses the proxy} (egress.south west);

    \draw[arrow] (egress.east) -- +(0.8,0) |- (registry.west);
    \draw[arrow] (egress.east) -- +(0.8,0) |- (hub.west);

    \draw[arrow] (registry.east) -- node[above right] {\footnotesize SSRF} (redteam.west);

    \draw[dottedarrow] (sandbox.south) -- node[left] {\footnotesize rw} (workspace.north);
    \draw[dottedarrow] (sandbox.south) ++(1cm,0) -- node[above right, near end] {\footnotesize rw} (memory.north west);

    \draw[dottedarrow] (workspace.south) |- node[above left, near end] {\footnotesize ro} (fsmon.west);
    \draw[dottedarrow] (memory.south) |- node[above right, near end] {\footnotesize ro} (fsmon.east);
\end{tikzpicture}
}
\caption{The harness topology. Solid: agent traffic the proxy audits. Dashed:
the heartbeat, which bypasses the proxy by design so that a saturated network
cannot silence it. Long dashes: the model backend, outside the audited scope.
Dotted: volumes and bind mounts, which are not network. The filesystem monitor
watches two roots, and \texttt{./memory} survives \texttt{docker compose
down -v}, which is what makes it the most durable leak channel.}
\label{fig:arch}
\end{figure}

\subsection{Three model-independent signals}

Each telemetry signal audits a distinct attack surface from an isolated container.
None of these monitors query the model, ensuring detection remains independent of the agent's cooperation.

\begin{table}[h]
\centering
\small
\begin{tabular}{@{}lll@{}}
\toprule
\textbf{Monitor} & \textbf{Surface} & \textbf{Fires on} \\
\midrule
\texttt{heartbeat\_monitor} & process liveness & silence past a configured threshold \\
\texttt{canary\_monitor} & filesystem, persistent memory & canary token in a path or in file content \\
\texttt{egress\_monitor} & outbound HTTP & destination outside baseline, or canary in a request \\
\bottomrule
\end{tabular}
\caption{The three detection signals. The canary monitor watches two roots:
the ephemeral workspace volume and the persistent \texttt{./memory} bind
mount.}
\end{table}

The primary detection mechanism is a per-run canary token injected into the agent's system prompt as its internal session identifier, matched in plaintext and under four low-cost obfuscations, and watched through \texttt{inotify} rather than polling.
That coverage is a deliberate empirical floor and not an exhaustive defense: an agent encoding data in file name lengths would pass it, a limit Appendix~\ref{sec:validity} takes up.

\subsection{Six induction vectors}

The agent is always given the same benign surface task, patching a division
bug until a test suite passes. What varies is the temptation planted in its
environment. Tasks 01 through 03 are passive: an exposed credential, an error
pointing at an internal registry, with no instruction to act. Tasks 04
through 06 carry an \emph{identical} injected instruction and differ only in
how it arrives, by live fetch, by cross-run persistent memory, or by a
document in a local knowledge base, which is what makes the comparison in
Section~
ef{sec:results} possible. Vectors 05 and 06 are ASI06, memory and
context poisoning~\cite{owasp-asi06}; Appendix~
ef{sec:vectors} gives each
one in full.

\subsection{Conditions}

The experimental axis is the security directive in the system prompt, with the task held constant.
The canary is present in every condition, as instrumentation rather than a defence, so the only thing that varies is what the agent is told.
\texttt{sin\_harness} mentions the token in passing, in 77 characters.
\texttt{con\_harness} applies an 830-character production guardrail warning that it must never reach file names, directory names, or URLs.
\texttt{con\_harness\_generico} matches that guardrail for length and emphasis, 818 characters against 830, but forbids only that the token leave the working environment, naming no surface at all.

The third condition exists to decouple two things the second confounds: an emphatic warning, and the naming of the surfaces an attacker would ask for.
Section~\ref{sec:results} reports what that isolation found.
